\begin{abstract}
Shift-start personal protective equipment (PPE) and fatigue screening in Northern Ontario mining remains manual, inconsistent, and undocumented. Existing computer vision systems detect PPE globally but fail to account for zone-specific requirements under Ontario Regulation~854---producing false alarms when workers enter surface areas where hard hats are not required.

We present \textbf{Northern Shift Guard}, a three-layer vision-language decision framework for explainable and auditable shift-start safety screening. Given a worker photo and selected mine zone, the system: (1)~extracts structured PPE and fatigue evidence via GPT-4o vision, (2)~checks detections against zone-specific compliance rules, and (3)~reasons over evidence using NVIDIA Nemotron to produce prioritized supervisor actions grounded in regulatory context. Every decision is stored in an audit log with full evidence chain.

We evaluate on \textbf{[N]} images across \textbf{[M]} mine zones and show that zone-aware compliance reduces false alarm rates by \textbf{[X]\%} compared to global-rule baselines, while maintaining \textbf{[Y]\%} PPE detection accuracy. A zone-flip case study demonstrates that identical visual evidence yields different compliance outcomes under surface versus active stope requirements. We discuss limitations including fatigue screening as a visual aid (not medical diagnosis), privacy considerations, and deployment requirements for production mine sites.
\end{abstract}

\textbf{Keywords:} mining safety, PPE compliance, explainable AI, vision-language model, zone-aware screening, audit trail, shift-start safety
